\begin{example} % 烂题
（2026朝阳二模）已知椭圆 $E:\frac{x^2}{a^2}+\frac{y^2}{b^2}=1\ (a>b>0)$ 的长轴长与短轴长之和为 $6$，焦距为 $2\sqrt{3}$。\\
（1）求椭圆 $E$ 的方程；\\
（2）设 $O$ 为原点，点 $A(x_1,y_1),B(x_2,y_2)$ 分别为椭圆 $E$ 上位于第一象限、第二象限内的点，且 $x_1=2y_2$。当点 $M$ 满足 $\overrightarrow{OM}=\frac{1}{2}\overrightarrow{OA}+\frac{\sqrt{3}}{2}\overrightarrow{OB}$ 时，求证：点 $M$ 在椭圆 $E$ 上。
\end{example}

\begin{example} % 动点引入，韦达定理求坐标，三点共线
    % 斜率例题
（2026东城二模）已知椭圆 $W:\frac{x^2}{a^2}+\frac{y^2}{b^2}=1\ (a>b>0)$ 的右焦点为 $(\sqrt{3},0)$，且经过点 $\left(-1,\frac{\sqrt{3}}{2}\right)$。\\
（1）求椭圆 $W$ 的方程；\\
（2）已知点 $A(a,0),T(0,t)\ (t\neq\pm1)$。设直线 $x=a$ 分别与直线 $y=b,y=-b$ 交于点 $M,N$，直线 $MT$ 与直线 $y=-b$ 交于点 $P$，直线 $NT$ 与直线 $y=b$ 交于点 $Q$，直线 $AT$ 交椭圆 $W$ 于另一点 $R$。求证：$P,Q,R$ 三点共线。
\end{example}

\begin{example} % 直线引入，
    % 定点例题
（2026丰台二模）已知椭圆 $E:\frac{x^2}{a^2}+\frac{y^2}{b^2}=1\ (a>b>0)$ 的离心率为 $\frac{\sqrt{3}}{2}$，以椭圆 $E$ 的焦点和短轴顶点为顶点的四边形是边长为 $2$ 的菱形。\\
（1）求椭圆 $E$ 的标准方程；\\
（2）已知 $A$ 为椭圆 $E$ 的左顶点，$M,N$ 为椭圆 $E$ 上两个不同的动点（均不与点 $A$ 重合），且直线 $AM$ 与直线 $AN$ 的斜率之积为 $\frac{5}{4}$。求证：直线 $MN$ 过定点。
\end{example}

\begin{example} % 直线或者动点引入，三角形面积的转换（千锤法）
% 入题思路例题1.9
% 条件翻译面积例题
（2026海淀二模）椭圆 $E:\frac{x^2}{a^2}+\frac{y^2}{b^2}=1\ (a>b>0)$ 的下顶点为 $A(0,-1)$，左焦点到点 $A$ 的距离为 $\sqrt{3}$。\\
（1）求椭圆 $E$ 的方程及离心率；\\
（2）经过点 $A$ 的直线与椭圆的另一个交点为 $D$（位于 $x$ 轴上方），与 $x$ 轴的交点为 $M$。点 $D$ 关于 $x$ 轴的对称点为 $E$。过点 $D$ 作与 $x$ 轴平行的直线交直线 $AE$ 于点 $N$。设$\triangle DEN$与$\triangle AEM$的面积分别为 $S_{\triangle DEN}$ 和 $S_{\triangle AEM}$，若 $S_{\triangle DEN}=3S_{\triangle AEM}$，求直线 $AD$ 的斜率。
\end{example}

\begin{example} % 直线引入，三角形面积的转换
    % 练习题目
（2026西城二模）已知椭圆 $C:\frac{x^2}{a^2}+\frac{y^2}{b^2}=1\ (a>b>0)$ 的左焦点为 $F(-1,0)$，点 $M\left(-1,\frac{2\sqrt{3}}{3}\right)$ 在椭圆 $C$ 上。\\
（1）求椭圆 $C$ 的方程；\\
（2）过点 $N(-3,0)$ 的直线与椭圆 $C$ 交于 $A,B$ 两点，过点 $A$ 作 $AP$垂直直线$ MF$ 于点 $P$。记 $\triangle BPN$ 和 $\triangle BPF$ 的面积分别为 $S_{\triangle BPN}$ 和 $S_{\triangle BPF}$，求证：$S_{\triangle BPN}=S_{\triangle BPF}$。
\end{example}

\begin{example} % 直线引入，倾斜角互余
    % 斜率例题
（2025朝阳二模）已知椭圆 $E:\frac{x^2}{a^2}+\frac{y^2}{b^2}=1\ (a>b>0)$ 的焦距为 $2$，且过点 $(0,\sqrt{3})$。\\
（1）求椭圆 $E$ 的方程；\\
（2）设直线 $y=kx+2\ (k<0)$ 与椭圆 $E$ 交于不同的两点 $A,B$，直线 $y=x$ 与直线 $AB$ 交于点 $N$。若 $\angle AON=\angle BON$（$O$ 为坐标原点），求 $k$ 的值。
\end{example}

\begin{example} % 直线引入，三角形面积
% 入题思路例题1.5
% 条件翻译面积例题
（2025东城二模）已知椭圆 $E:\frac{x^2}{a^2}+\frac{y^2}{b^2}=1\ (a>b>0)$ 的一个顶点为 $A(0,1)$，且过点 $\left(1,\frac{\sqrt{3}}{2}\right)$。\\
（1）求椭圆 $E$ 的方程及焦距；\\
（2）过点 $(0,2)$ 的直线与椭圆 $E$ 交于不同的两点 $B,C$。直线 $AB,AC$ 的斜率分别记为 $k_1,k_2$，当 $k_1+k_2=\frac{2}{3}$ 时，求 $\triangle ABC$ 的面积。
\end{example}

\begin{example} % 直角三角形特殊角
    % 特殊三角形例题
（2025海淀二模）已知椭圆 $C:\frac{x^2}{6}+\frac{y^2}{2}=1$。设直线 $l:y=x+m$ 交椭圆 $C$ 于不同的两点 $A,B$，与 $y$ 轴交于点 $P$。\\
（1）当 $m=0$ 时，求 $\lvert AB\rvert$ 的值；\\
（2）若点 $Q$ 满足 $\lvert PQ\rvert=3$ 且 $\lvert QA\rvert=\lvert QB\rvert$，求 $\angle AQB$ 的大小。
\end{example}

\begin{example} % 动点或者直线引入，三角形面积
% 入题思路例题1.10
% 练习题
（2025丰台二模）已知椭圆 $E:\frac{x^2}{a^2}+\frac{y^2}{b^2}=1\ (a>b>0)$ 的左顶点为 $A(-2,0)$，焦距为 $2\sqrt{2}$。\\
（1）求椭圆 $E$ 的方程；\\
（2）设 $O$ 为原点，过点 $A$ 且斜率为 $k$ 的直线 $l$ 与椭圆 $E$ 的另一个交点为 $T$，线段 $AT$ 的垂直平分线与 $x$ 轴交于点 $M$，与 $y$ 轴交于点 $N$。过点 $P(1,0)$ 且与 $l$ 平行的直线与 $y$ 轴交于点 $Q$。若 $\triangle OMN$ 与 $\triangle NPQ$ 的面积之比为 $4:3$，求 $k$ 的值。
\end{example}

\begin{example} % 直线，计算量大
（2025西城二模）已知椭圆 $E:\frac{x^2}{a^2}+\frac{y^2}{b^2}=1\ (a>b>0)$，直线 $x+\sqrt{2}y+2\sqrt{2}=0$ 经过椭圆 $E$ 的左顶点 $A$ 和下顶点 $B$。\\
（1）求椭圆 $E$ 的方程和离心率；\\
（2）设过点 $G(0,s)\ (s>0)$ 且斜率不为 $0$ 的直线交椭圆 $E$ 于 $C,D$ 两点，直线 $BC,BD$ 与直线 $y=t$ 的交点分别为 $P,Q$，线段 $CD,PQ$ 的中点分别为 $M,N$。若直线 $MN$ 经过坐标原点，求 $s+t$ 的取值范围。
\end{example}

\begin{example} % 直线引入，菱形的翻译
    % 平行四边形，定点例题
（2024海淀二模）已知椭圆 $E$ 的焦点在 $x$ 轴上，中心在坐标原点。以 $E$ 的一个顶点和两个焦点为顶点的三角形是等边三角形，且其周长为 $6\sqrt{2}$。\\
（1）求椭圆 $E$ 的方程；\\
（2）设过点 $M(2,0)$ 的直线 $l$（不与坐标轴垂直）与椭圆 $E$ 交于不同的两点 $A,C$，与直线 $x=16$ 交于点 $P$。点 $B$ 在 $y$ 轴上，$D$ 为坐标平面内的一点，四边形 $ABCD$ 是菱形。求证：直线 $PD$ 过定点。
\end{example}

\begin{example} % 直线引入，交线求坐标、三点共线
    % 练习题
（2024朝阳二模）已知椭圆 $E$ 的两个顶点分别为 $A(-2,0),B(2,0)$，焦点在 $x$ 轴上，且椭圆 $E$ 过点 $C\left(\sqrt{2},\frac{\sqrt{2}}{2}\right)$。\\
（1）求椭圆 $E$ 的方程；\\
（2）设 $O$ 为原点，不经过椭圆 $E$ 顶点的直线 $l$ 与椭圆 $E$ 交于两点 $P(x_1,y_1),Q(x_2,y_2)\ (x_1\neq x_2)$，直线 $BP$ 与直线 $OC$ 交于点 $H$，点 $M$ 与点 $Q$ 关于原点对称。\\
（i）求点 $H$ 的坐标（用 $x_1,y_1$ 表示）；\\
（ii）若 $A,H,M$ 三点共线，求证：直线 $l$ 经过定点。
\end{example}

\begin{example} % 烂题
（2024东城二模）已知椭圆 $C:\frac{x^2}{a^2}+\frac{y^2}{b^2}=1\ (a>b>0)$ 的右焦点为 $F(\sqrt{2},0)$，左、右顶点分别为 $A,B$，直线 $l:x=2\sqrt{2}$，且点 $A$ 到点 $F$ 的距离与点 $A$ 到直线 $l$ 的距离之比为 $\frac{\sqrt{2}}{2}$。\\
（1）求椭圆 $C$ 的方程；\\
（2）设 $M,N$ 为椭圆 $C$ 上不同的两点（不在坐标轴上），过点 $N$ 作直线 $BM$ 的平行线，与直线 $AM$ 交于点 $D$；过点 $M$ 作直线 $BN$ 的平行线，与直线 $AN$ 交于点 $E$。求证：点 $D$ 与点 $E$ 到直线 $l$ 的距离相等。
\end{example}

\begin{example} % 动点引入，韦达定理求坐标，三点共线
    % 练习题
（2024西城二模）已知椭圆 $E:\frac{x^2}{a^2}+\frac{y^2}{b^2}=1\ (a>b>0)$ 的一个顶点为 $D(2,0)$，焦距为 $2\sqrt{2}$。\\
（1）求椭圆 $E$ 的方程；\\
（2）设点 $P$ 是第一象限内椭圆 $E$ 上一点，过点 $P$ 作 $y$ 轴的垂线，垂足为 $Q$。点 $P$ 关于原点的对称点为 $A$，直线 $DQ$ 与椭圆的另一个交点为 $B$，直线 $DP$ 与 $y$ 轴的交点为 $C$。求证：$A,B,C$ 三点共线。
\end{example}

\begin{example} % 直线引入（题目有误导），中点的处理
    % 练习题目
（2024丰台二模）已知两点 $F_1(-1,0),F_2(1,0)$，曲线 $\Omega$ 上的动点 $M$ 满足 $\lvert MF_1\rvert+\lvert MF_2\rvert=2\lvert F_1F_2\rvert$，直线 $MF_2$ 与曲线 $\Omega$ 交于另一点 $N$。\\
（1）求曲线 $\Omega$ 的方程；\\
（2）设曲线 $\Omega$ 与 $x$ 轴的交点分别为 $A,B$（点 $A$ 在点 $B$ 的左侧，且 $M$ 不与 $A,B$ 重合），直线 $AM$ 与直线 $BN$ 交于点 $P$。当点 $B$ 为线段 $NP$ 的中点时，求点 $N$ 的横坐标。
\end{example}

\begin{example} % 直线引入，铅锤法求面积
% 条件翻译面积例题
（2023西城二模）已知椭圆 $E:\frac{x^2}{a^2}+\frac{y^2}{b^2}=1\ (a>b>0)$ 的短轴长为 $2\sqrt{2}$，一个焦点为 $F_1(-2,0)$。\\
（1）求椭圆 $E$ 的方程和离心率；\\
（2）设直线 $l:x-my-2=0$ 与椭圆 $E$ 交于两点 $A,B$，点 $M$ 在线段 $AB$ 上，点 $F_1$ 关于点 $M$ 的对称点为 $C$。当四边形 $AF_1BC$ 的面积最大时，求 $m$ 的值。
\end{example}

\begin{example} % 动点引入，交线求坐标，弦长公式
（2023朝阳二模）已知点 $\left(1,\frac{3}{2}\right)$ 在椭圆 $E:\frac{x^2}{a^2}+\frac{y^2}{b^2}=1\ (a>b>0)$ 上，且 $E$ 的离心率为 $\frac{1}{2}$。\\
（1）求椭圆 $E$ 的方程；\\
（2）设 $F$ 为椭圆 $E$ 的右焦点，点 $P(m,n)$ 是 $E$ 上的任意一点，直线 $PF$ 与直线 $3mx+4ny=0$ 相交于点 $Q$，求 $\lvert PQ\rvert$ 的值。
\end{example}

\begin{example} % 动点引入，角度的翻译
（2023海淀二模）已知椭圆 $E:\frac{x^2}{a^2}+\frac{y^2}{b^2}=1\ (a>b>0)$ 的左顶点为 $A$，上、下顶点分别为 $B_1,B_2$，直线 $AB_1$ 的方程为 $x-\sqrt{3}y+\sqrt{3}=0$。\\
（1）求椭圆 $E$ 的方程及离心率；\\
（2）$P$ 是椭圆上一点，且在第一象限内，$M$ 是点 $P$ 关于 $x$ 轴的对称点。过点 $P$ 作垂直于 $y$ 轴的直线，交直线 $AB_1$ 于点 $Q$；再过点 $Q$ 作垂直于 $x$ 轴的直线，交直线 $PB_2$ 于点 $N$。求 $\angle MNQ$ 的大小。
\end{example}

\begin{example} % 直线引入，三点共线，垂直条件
    % 练习题
（2023丰台二模）已知椭圆 $C:\frac{x^2}{a^2}+\frac{y^2}{b^2}=1\ (a>b>0)$ 经过两点 $A(-2,0),B\left(\sqrt{3},\frac{1}{2}\right)$。\\
（1）求椭圆 $C$ 的方程和离心率；\\
（2）设 $P,Q$ 为椭圆 $C$ 上不同的两个点，直线 $AP$ 与 $y$ 轴交于点 $E$，直线 $AQ$ 与 $y$ 轴交于点 $F$。若点 $M(1,0)$ 满足 $MF\perp ME$，求证：$P,O,Q$ 三点共线。
\end{example}

\begin{example} % 入题思路（我搞复杂了），直线过顶点
% 入题思路例题1.14，定点例题
（2022朝阳二模）已知椭圆 $C:\frac{x^2}{a^2}+\frac{y^2}{b^2}=1\ (a>b>0)$ 的一个顶点为 $P(0,1)$，离心率为 $\frac{\sqrt{2}}{2}$。\\
（1）求椭圆 $C$ 的方程；\\
（2）过点 $P$ 作斜率为 $k_1$ 的直线 $l_1$，交椭圆 $C$ 于另一点 $A$；过点 $P$ 作斜率为 $k_2\ (k_2\neq k_1)$ 的直线 $l_2$，交椭圆 $C$ 于另一点 $B$。若 $k_1k_2=1$，求证：直线 $AB$ 经过定点。
\end{example}

\begin{example} % 直线引入，平行四边形的翻译
    % 练习题
（2022海淀二模）椭圆 $M:\frac{x^2}{a^2}+\frac{y^2}{b^2}=1\ (a>b>0)$ 的左顶点为 $A(-2,0)$，离心率为 $\frac{\sqrt{3}}{2}$。\\
（1）求椭圆 $M$ 的方程；\\
（2）已知经过点 $\left(0,\frac{\sqrt{3}}{2}\right)$ 的直线 $l$ 交椭圆 $M$ 于 $B,C$ 两点，$D$ 是直线 $x=-4$ 上一点。若四边形 $ABCD$ 为平行四边形，求直线 $l$ 的方程。
\end{example}

\begin{example} % 动点引入，椭圆与圆的关系，基本不等式求最值
    % 圆例题
（2022西城二模）已知椭圆 $C:\frac{x^2}{a^2}+\frac{y^2}{b^2}=1\ (a>b>0)$ 的左顶点为 $A(-2,0)$，圆 $O:x^2+y^2=1$ 经过椭圆 $C$ 的上、下顶点。\\
（1）求椭圆 $C$ 的方程和焦距；\\
（2）已知 $P,Q$ 分别是椭圆 $C$ 和圆 $O$ 上的动点（$P,Q$ 不在坐标轴上），且直线 $PQ$ 与 $x$ 轴平行，线段 $AP$ 的垂直平分线与 $y$ 轴交于点 $M$，圆 $O$ 在点 $Q$ 处的切线与 $y$ 轴交于点 $N$。求线段 $MN$ 长度的最小值。
\end{example}

\begin{example} %直线引入，倾斜角互余
    % 练习题
（2022东城二模）已知椭圆 $E:\frac{x^2}{a^2}+\frac{y^2}{b^2}=1\ (a>b>0)$ 的右顶点为 $A(2,0)$，离心率为 $\frac{1}{2}$。过点 $P(6,0)$ 且与 $x$ 轴不重合的直线 $l$ 交椭圆 $E$ 于不同的两点 $B,C$，直线 $AB,AC$ 分别交直线 $x=6$ 于点 $M,N$。\\
（1）求椭圆 $E$ 的方程；\\
（2）设 $O$ 为原点，求证：$\angle PAN+\angle POM=90^\circ$。
\end{example}

\begin{example} % 直线引入，三点共线，直线消元
（2022丰台二模）已知椭圆 $C:\frac{x^2}{a^2}+\frac{y^2}{b^2}=1\ (a>b>0)$ 经过点 $P(2,1)$，点 $P$ 到椭圆 $C$ 的两个焦点的距离和为 $4\sqrt{2}$。\\
（1）求椭圆 $C$ 的方程；\\
（2）设 $Q(4,0)$，$R$ 为 $PQ$ 的中点，作 $PQ$ 的平行线 $l$，与椭圆 $C$ 交于不同的两点 $A,B$。直线 $AQ$ 与椭圆 $C$ 交于另一点 $M$，直线 $BQ$ 与椭圆 $C$ 交于另一点 $N$。求证：$M,N,R$ 三点共线。
\end{example}

\begin{example} % 动点引入，三角形面积比（同底不同高）
    % 练习题目
（2021西城二模）已知椭圆 $C:\frac{x^2}{a^2}+\frac{y^2}{b^2}=1$ 的离心率为 $\frac{\sqrt{6}}{3}$，其长轴的两个端点分别为 $A(-3,0),B(3,0)$。\\
（1）求椭圆 $C$ 的标准方程；\\
（2）点 $P$ 为椭圆上除 $A,B$ 外的任意一点，直线 $AP$ 交直线 $x=4$ 于点 $E$，点 $O$ 为坐标原点。过点 $O$ 且与直线 $BE$ 垂直的直线记为 $l$，直线 $BP$ 交 $y$ 轴于点 $M$，交直线 $l$ 于点 $N$。求 $\triangle BMO$ 与 $\triangle NMO$ 的面积之比。
\end{example}

\begin{example} % 直线引入，定值定点问题
（2021东城二模）已知椭圆 $C:\frac{x^2}{a^2}+\frac{y^2}{b^2}=1\ (a>b>0)$ 的右焦点为 $F$，左、右顶点分别为 $A(-2,0),B(2,0)$，$\lvert AF\rvert=3\lvert FB\rvert$。\\
（1）求椭圆 $C$ 的方程；\\
（2）设过点 $P(2,1)$ 的直线 $l$ 与椭圆 $C$ 交于不同的两点 $M,N$，过点 $N$ 作 $x$ 轴的垂线，与直线 $BM$ 交于点 $D$，$E$ 为线段 $DN$ 的中点。证明：直线 $BE$ 的斜率为定值。
\end{example}

\begin{example} % 直线引入，三点共线
（2021朝阳二模）已知 $F$ 为椭圆 $C:\frac{x^2}{2}+y^2=1$ 的左焦点，直线 $l:y=k(x-2)$ 与椭圆 $C$ 交于不同的两点 $M,N$。\\
（1）当 $k=-\frac{1}{2}$ 时，求 $\triangle FMN$ 的面积；\\
（2）设直线 $FM,FN$ 分别与直线 $x=1$ 交于两点 $P,Q$，线段 $MN,PQ$ 的中点分别为 $G,H$，点 $A\left(\frac{1}{5},0\right)$。当 $k$ 变化时，证明：$A,G,H$ 三点共线。
\end{example}

\begin{example} % 动点引入（入题思路比较难），垂直的翻译，弦长公式
% 入题思路例题1.12
% 距离公式例题
（2021海淀二模）椭圆 $C:\frac{x^2}{a^2}+\frac{y^2}{b^2}=1\ (a>b>0)$ 的左、右焦点分别为 $F_1,F_2$，$E$ 是椭圆 $C$ 上一点，且 $\lvert F_1F_2\rvert=2$，$\lvert EF_1\rvert+\lvert EF_2\rvert=4$。\\
（1）求椭圆 $C$ 的方程；\\
（2）$M,N$ 是 $y$ 轴上的两个动点（点 $M$ 与点 $E$ 位于 $x$ 轴的两侧），$\angle MF_1N=\angle MEN=90^\circ$，直线 $EM$ 交 $x$ 轴于点 $P$，求 $\frac{\lvert EP\rvert}{\lvert PM\rvert}$ 的值。
\end{example}

\begin{example} % 直线引入，定值定点（两个未知数的恒等式）
（2021丰台二模）已知椭圆 $C:\frac{x^2}{3}+y^2=1$，过点 $(-1,0)$ 的直线 $l$ 交椭圆 $C$ 于点 $A,B$。\\
（1）当直线 $l$ 与 $x$ 轴垂直时，求 $\lvert AB\rvert$；\\
（2）在 $x$ 轴上是否存在定点 $P$，使 $\overrightarrow{PA}\cdot\overrightarrow{PB}$ 为定值？若存在，求点 $P$ 的坐标及 $\overrightarrow{PA}\cdot\overrightarrow{PB}$ 的值；若不存在，说明理由。
\end{example}
